% !TEX root = ../main.tex
%% Chapitre 16 — Dénombrement
%% Programme officiel MPSI/MP2I (BO spécial n°1 du 11/02/2021), deuxième semestre,
%% section « Dénombrement » : a) cardinal d'un ensemble fini (parties, applications
%% entre ensembles de même cardinal, opérations sur les cardinaux, applications
%% d'un ensemble fini dans un autre, parties d'un ensemble fini) ; b) listes et
%% combinaisons (p-listes, arrangements, permutations, combinaisons,
%% démonstrations combinatoires de Pascal et du binôme).
%% La formule du crible générale est hors programme (cas de 2 et 3 ensembles en
%% exercice) ; les bijections ne sont pas un attendu (utilisées en approfondissement).
%% Résultats vérifiés par verifications/chapitre_16.py.

\begin{ouverture}[denombrement]{16}{Dénombrement}{Partie 2 · Deuxième semestre}
  {Cardinaux\enspace·\enspace Principe multiplicatif\enspace·\enspace Listes\enspace·\enspace
   Arrangements\enspace·\enspace Combinaisons\enspace·\enspace Binôme}
  \objectif{Manier les cardinaux : unions, complémentaires, produits cartésiens.}
  \objectif{Compter les applications, les injections et les parties d'un ensemble fini.}
  \objectif{Distinguer listes, arrangements, permutations et combinaisons — et choisir le bon modèle.}
  \objectif{Démontrer combinatoirement la formule de Pascal et celle du binôme.}
  \objectif{Pratiquer le double comptage et le passage au complémentaire.}
  \objectif{Préparer les probabilités finies du chapitre 17.}
  \prerequis{Ensembles, applications, injections/surjections/bijections (chapitre 1).}
  \prerequis{Coefficients binomiaux et formule du binôme, côté calcul (chapitre 2).}
  \prerequis{Récurrences (chapitres 1 et 5) ; le groupe symétrique (chapitre 14).}
\end{ouverture}

%% =====================================================================
%%  L'ESSENTIEL
%% =====================================================================
\essentiel

\essSection{1}{Cardinal d'un ensemble fini}

\begin{definition}[Cardinal]
Un ensemble $E$ est \textbf{fini} s'il est vide ou en bijection avec un
$\ens{1,\dots,n}$ ; l'entier $n$, unique, est son \textbf{cardinal}, noté
$|E|$ ou $\operatorname{Card}(E)$ (tout fondement théorique est hors
programme : on compte, et les propriétés intuitives sont admises). Toute
partie $A$ d'un ensemble fini $E$ est finie, avec $|A|\leqslant|E|$, et
\emph{l'égalité des cardinaux force l'égalité des ensembles} :
$|A|=|E|\implies A=E$.
\end{definition}

\begin{theoreme}[Applications entre ensembles de même cardinal]
Si $|E|=|F|$ est fini, une application $f:E\to F$ est bijective si et
seulement si elle est injective, si et seulement si elle est surjective —
le pendant ensembliste du théorème du rang (chapitre 12). Conséquence
immédiate, le \emph{principe des tiroirs} : si l'on range $n+1$ objets
dans $n$ tiroirs, l'un des tiroirs reçoit au moins deux objets
(exercice 9).
\end{theoreme}

\essSection{2}{Les opérations sur les cardinaux}

\noindent
\begin{minipage}[c]{0.55\linewidth}
\begin{propriete}[Sommes, différences, produits]
Pour $A,B$ parties d'un ensemble fini $E$ :
\begin{itemize}[label={\color{otNavy}\textbullet},leftmargin=1.4em,itemsep=0.6mm,topsep=0.8mm]
  \item union \emph{disjointe} : $|A\cup B|=|A|+|B|$ ; plus généralement
    une partition se somme ;
  \item union quelconque :
    $|A\cup B|=|A|+|B|-|A\cap B|$ (on a compté $A\cap B$ deux fois) ;
  \item complémentaire : $|\overline A|=|E|-|A|$, et différence :
    $|A\setminus B|=|A|-|A\cap B|$ ;
  \item produit cartésien : $|A\times B|=|A|\times|B|$ — c'est le
    \emph{principe multiplicatif} : une suite de choix indépendants se
    compte en multipliant les nombres d'options.
\end{itemize}
\end{propriete}
\end{minipage}\hfill
\begin{minipage}[c]{0.41\linewidth}
\centering
\begin{tikzpicture}[scale=0.88,line cap=round]
  \fill[otInk] (0,1.5) circle (1.6pt);
  \foreach \y/\lab in {2.6/e, 1.5/p, 0.4/c}{%
    \fill[otNavy] (1.35,\y) circle (1.6pt);
    \draw[otNavy,line width=0.7pt] (0,1.5) -- (1.35,\y);
    \foreach \d in {0.3,-0.3}{%
      \fill[otGreenDark] (2.7,{\y+\d}) circle (1.4pt);
      \draw[otGreenDark,line width=0.6pt] (1.35,\y) -- (2.7,{\y+\d});}}
  \node[font=\scriptsize,text=otNavy] at (1.35,3.0) {$3$ plats};
  \node[font=\scriptsize,text=otGreenDark] at (2.78,3.3) {$2$ desserts};
  \node[font=\scriptsize,text=otInk,anchor=east] at (-0.12,1.5) {départ};
  \node[font=\scriptsize,text=otMuted] at (1.5,-0.45) {$3\times2=6$ menus};
\end{tikzpicture}\\[0.5mm]
\legende{L'arbre des choix : chaque étage multiplie. C'est le principe
multiplicatif, cœur de tout le chapitre.}
\end{minipage}

\essSection{3}{Applications et parties d'un ensemble fini}

\begin{theoreme}[Deux comptes fondamentaux]
Soient $E$ et $F$ finis, $|E|=p$ et $|F|=n$.
\begin{itemize}[label={\color{otNavy}\textbullet},leftmargin=1.4em,itemsep=0.6mm,topsep=0.8mm]
  \item Le nombre d'applications de $E$ dans $F$ est
    $n^{p}$ : pour chacun des $p$ éléments, $n$ images possibles.
  \item Le nombre de parties de $E$ est
    $|\parties{E}|=2^{p}$ : une partie, c'est une application de $E$ dans
    $\ens{0,1}$ (sa fonction indicatrice, chapitre 1) — garder ou jeter
    chaque élément.
\end{itemize}
\end{theoreme}

\essSection{4}{Listes, arrangements, permutations}

\begin{definition}[Avec ou sans répétition, avec ordre]
Une \textbf{$p$-liste} ($p$-uplet) d'éléments de $F$ est un élément de
$F^{p}$ : il y en a $n^{p}$. Une $p$-liste d'éléments \emph{deux à deux
distincts} (un \emph{arrangement}) correspond à une application
\emph{injective} de $\ens{1,\dots,p}$ dans $F$ : il y en a
\[
n\,(n-1)\cdots(n-p+1)=\frac{n!}{(n-p)!}
\qquad(0\ \text{si }p>n) :
\]
$n$ choix, puis $n-1$, etc. Pour $p=n$, on retrouve les $n!$
\textbf{permutations} de $F$ — le cardinal de $S_n$ (chapitre 14).
\end{definition}

\begin{exemple}[Anagrammes]
Les anagrammes de MPSI : $4!=24$. Celles d'OPTITAB : $7$ lettres dont deux
T indiscernables, soit $\dfrac{7!}{2!}=2520$ — on divise par les
permutations internes des lettres répétées.
\end{exemple}

\essSection{5}{Combinaisons et binôme}

\begin{definition}[Combinaisons]
Une \textbf{combinaison} de $p$ éléments de $E$ ($|E|=n$) est une
\emph{partie} à $p$ éléments : ni ordre, ni répétition. Leur nombre est
$\dbinom np$, et comme chaque partie à $p$ éléments s'ordonne de $p!$
façons pour donner les arrangements :
\[
\binom np=\frac{n\,(n-1)\cdots(n-p+1)}{p!}=\frac{n!}{p!\,(n-p)!},
\qquad
\binom np=\binom{n}{n-p}.
\]
\end{definition}

\begin{demonstration}[Pascal et le binôme, preuves combinatoires]
\emph{Formule de Pascal.} Fixons un élément $a\in E$ avec $|E|=n$. Les
parties à $p$ éléments se partagent en celles qui contiennent $a$
(choisir les $p-1$ autres : $\binom{n-1}{p-1}$) et celles qui ne le
contiennent pas ($\binom{n-1}{p}$) :
$\binom np=\binom{n-1}{p-1}+\binom{n-1}{p}$.

\emph{Formule du binôme.} En développant $(a+b)^{n}=(a+b)\cdots(a+b)$,
chaque terme s'obtient en choisissant $a$ dans certains facteurs et $b$
dans les autres : le terme $a^{k}b^{\,n-k}$ apparaît autant de fois qu'il
y a de choix de $k$ facteurs parmi $n$, soit $\binom nk$ :
$(a+b)^{n}=\sum_{k=0}^{n}\binom nk a^{k}b^{\,n-k}$.
\end{demonstration}

\begin{exemple}
$\displaystyle\sum_{k=0}^{n}\binom nk=2^{n}$ (compter toutes les parties
par cardinal — ou $a=b=1$) et
$\displaystyle\sum_{k=0}^{n}(-1)^{k}\binom nk=0$ pour $n\geqslant1$
($a=-1$, $b=1$) : autant de parties de cardinal pair que de cardinal
impair (exercice 8).
\end{exemple}

\essSection{6}{La méthode : quel modèle ?}

\begin{savoirfaire}[Choisir le bon modèle de comptage]
Deux questions décident de tout : \emph{l'ordre compte-t-il ?} et
\emph{les répétitions sont-elles permises ?}
\begin{center}
\begin{tabularx}{0.92\linewidth}{@{}>{\otSemi\fontsize{8}{10}\selectfont\leavevmode\color{otNavy}}p{34mm}*{2}{>{\centering\arraybackslash}X}@{}}
 & \textbf{avec répétition} & \textbf{sans répétition}\\[0.6mm]
avec ordre & $n^{p}$ \ (listes) & $\dfrac{n!}{(n-p)!}$ \ (arrangements)\\[2.2mm]
sans ordre & \emph{(exercices 16, 17)} & $\dbinom np$ \ (combinaisons)\\
\end{tabularx}
\end{center}
Ensuite : passer au \emph{complémentaire} quand « au moins » apparaît ;
\emph{partitionner} selon un paramètre naturel ; et compter une même
quantité de deux façons — le \emph{double comptage} — pour démontrer une
identité (exercices 13, 14, 22).
\end{savoirfaire}

\begin{attention}
Ne jamais multiplier des choix qui se recouvrent : « choisir 4 personnes
dont au moins une femme » ne se compte pas en « une femme, puis 3
quelconques » (chaque comité à plusieurs femmes serait compté plusieurs
fois). Le bon réflexe : complémentaire, ou partition par le nombre exact
de femmes (exercice 10).
\end{attention}

%% =====================================================================
%%  EXERCICES
%% =====================================================================
\exercices
\rubrique{Applications}

\begin{exercice}[1]{Premiers décomptes}
\begin{questions}
  \item Un menu se compose d'une entrée parmi $3$, d'un plat parmi $4$ et
    d'un dessert parmi $2$ : combien de menus ?
  \item Combien de codes PIN à $4$ chiffres ?
  \item Combien de plaques d'immatriculation de la forme
    deux lettres, trois chiffres, deux lettres ?
\end{questions}
\end{exercice}

\begin{exercice}[1]{Unions et complémentaires}
Dans une classe de $30$ élèves, $18$ suivent l'option maths, $15$ l'option
physique, et $8$ les deux.
\begin{questions}
  \item Combien suivent au moins une option ? Aucune ?
  \item Énoncer et justifier la formule $|A\cup B|=|A|+|B|-|A\cap B|$.
\end{questions}
\end{exercice}

\begin{exercice}[1]{Listes, avec ou sans répétition}
\begin{questions}
  \item Combien de mots de $3$ lettres (alphabet de $26$) ? De mots de
    $3$ lettres deux à deux distinctes ?
  \item Combien d'applications de $\ens{1,2,3}$ dans $\ens{1,\dots,5}$ ?
    Combien d'injectives ?
\end{questions}
\end{exercice}

\begin{exercice}[1]{Permutations et anagrammes}
\begin{questions}
  \item Combien d'anagrammes du mot MPSI ? De façons de ranger $5$ livres
    sur une étagère ?
  \item Combien d'anagrammes du mot OPTITAB ? \emph{Attention aux deux T.}
\end{questions}
\end{exercice}

\begin{exercice}[1]{Combinaisons : premiers calculs}
\begin{questions}
  \item Calculer $\dbinom52$, $\dbinom{10}3$ et $\dbinom{12}4$.
  \item Combien de mains de $5$ cartes dans un jeu de $32$ ?
  \item Pourquoi l'ordre ne compte-t-il pas ici, et comment la formule en
    tient-elle compte ?
\end{questions}
\end{exercice}

\begin{exercice}[1]{Le triangle de Pascal}
\begin{questions}
  \item Construire le triangle de Pascal jusqu'à la ligne $n=6$.
  \item Y lire $\dbinom63$, la symétrie, et vérifier que la somme de la
    ligne $n$ vaut $2^{n}$.
  \item Redémontrer la formule de Pascal par un argument de comptage.
\end{questions}
\end{exercice}

\begin{exercice}[1]{Le binôme en action}
\begin{questions}
  \item Développer $(a+b)^{5}$, et donner le coefficient de $x^{3}$ dans
    $(1+x)^{7}$.
  \item Que donnent $a=b=1$, puis $a=-1$, $b=1$, dans la formule du
    binôme ?
\end{questions}
\end{exercice}

\begin{exercice}[1]{Compter des parties}
Soit $E$ de cardinal $n\geqslant1$.
\begin{questions}
  \item Retrouver $|\parties{E}|=2^{n}$ par le codage en indicatrices.
  \item Combien de parties contiennent un élément fixé $a$ ?
  \item Montrer qu'il y a autant de parties de cardinal pair que de
    cardinal impair.
\end{questions}
\end{exercice}

\rubrique{Entraînement}

\begin{exercice}[2]{Le principe des tiroirs}
\begin{questions}
  \item Montrer que parmi $13$ personnes, deux au moins sont nées le même
    mois.
  \item Montrer que parmi $n+1$ entiers choisis dans $\ens{1,\dots,2n}$,
    deux au moins sont consécutifs — donc premiers entre eux
    (chapitre 7).
  \item Montrer que parmi $5$ points d'un carré de côté $1$, deux au
    moins sont à distance au plus $\dfrac{\sqrt2}{2}$.
\end{questions}
\end{exercice}

\begin{exercice}[2]{Comités et délégations}
Un groupe compte $7$ femmes et $5$ hommes.
\begin{questions}
  \item Combien de comités de $4$ personnes avec exactement $2$ femmes ?
  \item Combien avec au moins une femme ?
  \item Compter de deux façons les comités de $4$ personnes munis d'un
    président, et en déduire une identité.
\end{questions}
\end{exercice}

\begin{exercice}[2]{Mains de cartes}
Dans un jeu de $32$ cartes ($4$ as) :
\begin{questions}
  \item Combien de mains de $5$ cartes avec exactement $2$ as ?
  \item Combien avec au moins un as ?
\end{questions}
\end{exercice}

\begin{exercice}[2]{Chemins sur une grille}
On va de $(0,0)$ à $(p,q)$ par pas « droite » ou « haut ».
\begin{questions}
  \item Montrer que le nombre de chemins est $\dbinom{p+q}{p}$.
  \item Combien de chemins de $(0,0)$ à $(4,3)$ ? Combien passent par
    $(1,2)$ ?
  \item Retrouver la formule de Pascal en regardant le dernier pas.
\end{questions}
\end{exercice}

\begin{exercice}[2]{Le comité et son président}
\begin{questions}
  \item Démontrer combinatoirement l'identité
    $k\dbinom nk=n\dbinom{n-1}{k-1}$.
  \item En déduire $\displaystyle\sum_{k=0}^{n}k\binom nk=n\,2^{n-1}$,
    puis $\displaystyle\sum_{k=0}^{n}k(k-1)\binom nk=n(n-1)\,2^{n-2}$.
\end{questions}
\end{exercice}

\begin{exercice}[2]{La formule de Vandermonde}
\begin{questions}
  \item Démontrer combinatoirement
    $\displaystyle\sum_{k=0}^{n}\binom pk\binom{q}{n-k}=\binom{p+q}{n}$
    — choisir $n$ personnes dans un groupe de $p$ femmes et $q$ hommes.
  \item En déduire
    $\displaystyle\sum_{k=0}^{n}\binom nk^{2}=\binom{2n}{n}$ : le
    coefficient binomial central des chapitres 11 et 15.
\end{questions}
\end{exercice}

\begin{exercice}[2]{Surjections sur deux ou trois éléments}
\begin{questions}
  \item Combien de surjections de $\ens{1,\dots,n}$ sur $\ens{1,2}$
    ($n\geqslant2$) ?
  \item Montrer que le nombre de surjections de $\ens{1,\dots,n}$ sur
    $\ens{1,2,3}$ vaut $3^{n}-3\cdot2^{n}+3$ ($n\geqslant3$).
    \emph{Retirer les applications qui évitent une valeur, en comptant
    bien les recouvrements.}
\end{questions}
\end{exercice}

\begin{exercice}[2]{Étoiles et barres}
\begin{questions}
  \item Montrer que le nombre de solutions entières positives ou nulles
    de $x_1+x_2+\dots+x_p=n$ est $\dbinom{n+p-1}{p-1}$. \emph{Coder une
    solution par une suite de $n$ étoiles et $p-1$ barres.}
  \item Combien de solutions pour $x+y+z=10$ ?
\end{questions}
\end{exercice}

\begin{exercice}[2]{Applications monotones}
Soient $p\leqslant n$.
\begin{questions}
  \item Montrer que le nombre d'applications \emph{strictement
    croissantes} de $\ens{1,\dots,p}$ dans $\ens{1,\dots,n}$ est
    $\dbinom np$.
  \item Montrer que le nombre d'applications \emph{croissantes} (au sens
    large) est $\dbinom{n+p-1}{p}$. \emph{Se ramener à l'exercice 16.}
\end{questions}
\end{exercice}

\begin{exercice}[2]{Le binôme rencontre l'analyse}
\begin{questions}
  \item En intégrant $(1+x)^{n}$ sur $\intFF{0}{1}$ (chapitre 15),
    montrer que
    $\displaystyle\sum_{k=0}^{n}\frac{1}{k+1}\binom nk
    =\frac{2^{n+1}-1}{n+1}$.
  \item En dérivant $(1+x)^{n}$, retrouver
    $\displaystyle\sum_{k=0}^{n}k\binom nk=n\,2^{n-1}$ (exercice 13).
\end{questions}
\end{exercice}

\rubrique{Approfondissement}

\begin{exercice}[3]{Les dérangements}
Un \emph{dérangement} est une permutation sans point fixe ; on note $D_n$
leur nombre dans $S_n$ ($D_0=1$, $D_1=0$).
\begin{questions}
  \item En classant les dérangements selon l'image de $n$ et le sort de
    son antécédent, établir
    $D_n=(n-1)\bigl(D_{n-1}+D_{n-2}\bigr)$ pour $n\geqslant2$.
  \item En déduire par récurrence $D_n=n\,D_{n-1}+(-1)^{n}$, puis
    \[
    D_n=n!\,\sum_{k=0}^{n}\frac{(-1)^{k}}{k!} .
    \]
  \item Montrer que $\dfrac{D_n}{n!}\to\dfrac{1}{\mathrm{e}}$
    (chapitres 5 et 15) : environ $37\,\%$ des permutations sont des
    dérangements.
\end{questions}
\end{exercice}

\begin{exercice}[3]{Le crible pour trois ensembles}
\begin{questions}
  \item À partir de la formule pour deux ensembles, démontrer
    \[
    |A\cup B\cup C|=|A|+|B|+|C|-|A\cap B|-|A\cap C|-|B\cap C|
    +|A\cap B\cap C| .
    \]
  \item Combien d'entiers de $1$ à $100$ ne sont divisibles ni par $2$,
    ni par $3$, ni par $5$ ? (Chapitre 7.)
\end{questions}
\end{exercice}

\begin{exercice}[3]{Chemins sous la diagonale : les nombres de Catalan}
On note $C_n$ le nombre de chemins de $(0,0)$ à $(n,n)$ (pas droite/haut)
qui ne montent jamais strictement au-dessus de la diagonale.
\begin{questions}
  \item Calculer $C_1,C_2,C_3$ à la main.
  \item Par l'argument de réflexion (mettre en bijection les chemins
    \emph{fautifs} avec les chemins de $(0,0)$ à $(n-1,n+1)$), montrer
    que $C_n=\dbinom{2n}{n}-\dbinom{2n}{n-1}=\dfrac{1}{n+1}\dbinom{2n}{n}$.
  \item Vérifier : $1,2,5,14,42,\dots$
\end{questions}
\end{exercice}

\begin{exercice}[3]{Identités par double comptage}
\begin{questions}
  \item Démontrer combinatoirement
    $\dbinom{2n}{2}=2\dbinom n2+n^{2}$ — couper une population de $2n$
    personnes en deux moitiés.
  \item Démontrer la « crosse de hockey »
    $\displaystyle\sum_{k=2}^{n}\binom k2=\binom{n+1}{3}$ — classer les
    parties à $3$ éléments de $\ens{1,\dots,n+1}$ selon leur maximum.
\end{questions}
\end{exercice}

\begin{exercice}[3]{Permutations et points fixes}
\begin{questions}
  \item Montrer que le nombre de permutations de $S_n$ ayant
    \emph{exactement} $k$ points fixes est $\dbinom nk D_{n-k}$
    (exercice 19).
  \item En déduire
    $\displaystyle\sum_{k=0}^{n}\binom nk D_{n-k}=n!$ — puis, en comptant
    les couples (permutation, point fixe) de deux façons, que le nombre
    \emph{moyen} de points fixes d'une permutation vaut $1$
    (un avant-goût du chapitre 17).
\end{questions}
\end{exercice}

\begin{exercice}[3]{Fibonacci compte aussi}
Soit $a_n$ le nombre de parties de $\ens{1,\dots,n}$ sans deux éléments
consécutifs ($\varnothing$ comprise).
\begin{questions}
  \item Calculer $a_1,a_2,a_3$, puis montrer que
    $a_n=a_{n-1}+a_{n-2}$ — selon que $n$ appartient ou non à la partie.
    En déduire $a_n=F_{n+2}$ (chapitres 5, 12 et 13).
  \item Montrer que le nombre de telles parties à $k$ éléments est
    $\dbinom{n-k+1}{k}$, et en déduire l'identité
    $\displaystyle\sum_{k\geqslant0}\binom{n-k+1}{k}=F_{n+2}$ : Fibonacci
    dans le triangle de Pascal, le long des diagonales.
\end{questions}
\end{exercice}

% !TEX root = ../../main.tex
%% Défi du chapitre 16 — Les nombres de Stirling de seconde espèce
%% Résultats vérifiés par verifications/defis_16_20.py.

\rubrique{Défi}

\begin{defi}{Les nombres de Stirling de seconde espèce}{1 h 30}
Pour $n\geqslant1$ et $k\geqslant1$, on note $S(n,k)$ le nombre de façons de
partager l'ensemble $\intFF{1}{n}$ en $k$ parties non vides (des
\emph{blocs}), sans tenir compte de l'ordre des blocs ; on convient que
$S(0,0)=1$, et $S(n,k)=0$ si $k>n$ ou si $k=0<n$. Ainsi $S(3,2)=3$ :
$\ens{1}\ens{2,3}$, $\ens{2}\ens{1,3}$, $\ens{3}\ens{1,2}$.

\partieExo{A}{Premières valeurs et récurrence}
\begin{questions}
  \item Déterminer $S(n,1)$, $S(n,n)$, puis montrer que
    $S(n,n-1)=\dbinom n2$ et $S(n,2)=2^{n-1}-1$.
  \item En distinguant si l'élément $n$ forme un bloc à lui seul ou non,
    établir la relation
    \[
    S(n,k)=S(n-1,k-1)+k\,S(n-1,k)\qquad(n,k\geqslant1).
    \]
    Retrouver ainsi les lignes $n=5$ et $n=6$ du triangle de la figure.
  \item Dessiner les $S(4,2)$ partitions de $\intFF{1}{4}$ en deux blocs.
\end{questions}

\medskip
\begin{center}
\begin{tikzpicture}[x=1.15cm,y=0.62cm]
  \foreach \row [count=\n] in {{1},{1,1},{1,3,1},{1,7,6,1},{1,15,25,10,1},{1,31,90,65,15,1}}{%
    \node[font=\scriptsize,text=otMuted,anchor=east] at (-0.9,-\n) {$n=\n$};
    \foreach \v [count=\k] in \row {%
      \node[font=\small,text=otInk,inner sep=1.5pt] (s\n\k) at ({\k-1-(\n-1)/2+2.5},-\n) {$\v$};}}
  \draw[otGreenDark,line width=0.8pt,rounded corners=2pt] ($(s53)+(-0.28,-0.32)$) rectangle ($(s53)+(0.28,0.32)$);
  \draw[otNavy,line width=0.8pt,-{Stealth[length=1.6mm]}] (s42) -- node[left,font=\tiny,text=otNavy]{$\times1$} (s53);
  \draw[otNavy,line width=0.8pt,-{Stealth[length=1.6mm]}] (s43) -- node[right,font=\tiny,text=otNavy]{$\times3$} (s53);
\end{tikzpicture}\\[0.5mm]
\legende{Le triangle des $S(n,k)$ ($k$ croissant de gauche à droite) :
$S(5,3)=S(4,2)+3\,S(4,3)=7+3\times6=25$.}
\end{center}

\partieExo{B}{Surjections et factorielles descendantes}
\begin{questions}[resume]
  \item Montrer que le nombre de surjections de $\intFF{1}{n}$ sur
    $\intFF{1}{k}$ vaut $k!\,S(n,k)$.
  \item En classant les applications de $\intFF{1}{n}$ dans $\intFF{1}{x}$
    ($x\in\N$) selon leur image, montrer que
    \[
    x^{n}=\sum_{j=0}^{n}S(n,j)\,x(x-1)\cdots(x-j+1) .
    \]
  \item En déduire l'égalité de polynômes
    $X^{n}=\sum_{j=0}^{n}S(n,j)\,X(X-1)\cdots(X-j+1)$, et que les
    coordonnées de $X^{n}$ dans la base binomiale $(H_j)$ du défi du
    chapitre 12 sont les $j!\,S(n,j)$. Vérifier sur $X^{3}$.
\end{questions}

\partieExo{C}{Nombres de Bell et formule explicite}
\begin{questions}[resume]
  \item On note $B_n=\sum_kS(n,k)$ le nombre total de partitions de
    $\intFF{1}{n}$ ($B_0=1$). En considérant le bloc qui contient $n+1$,
    montrer que $B_{n+1}=\sum_{k=0}^{n}\binom nkB_k$, et calculer
    $B_1,\dots,B_6$.
  \item À l'aide de l'exercice 15 (surjections sur trois éléments), montrer
    que $S(n,3)=\dfrac{3^{n}-3\cdot2^{n}+3}{6}$ et le vérifier sur le
    triangle.
\end{questions}
\end{defi}

